\section{Drugie twierdzenie o niezupełności}\label{sec:drugie-twierdzenie}
\begin{definicja}[Formalna niesprzeczność]\label{res:12:1}
Ustalmy $\bot:=(0=S(0))$, którego negacja jest dowodliwa w $Q$. \emph{Zdaniem niesprzeczności} teorii $T$, dla przyjętego standardowego predykatu dowodu, nazywamy $\operatorname{Con}_T:=\neg\operatorname{Bew}_T(\ulcorner\bot\urcorner)$, czyli $\neg\Box\bot$. W $\mathbb N$ jest ono prawdziwe dokładnie wtedy, gdy $T$ jest niesprzeczna: sprzeczność daje dowód $\bot$ przez zasadę eksplozji, a dowód $\bot$ wraz z $T\vdash\neg\bot$ daje sprzeczność. Jest to zdanie $\Pi_1$, czyli negacja zdania $\Sigma_1$.
\end{definicja}

\begin{twierdzenie}[Drugie twierdzenie Gödla]\label{res:12:2}
Niech $T\supseteq\mathrm{PA}$ będzie niesprzeczna i efektywnie aksjomatyzowalna. Dla standardowego predykatu dowodliwości, spełniającego (D1)--(D3) na mocy Lematu~\ref{res:11:6}, zachodzi $T\nvdash\operatorname{Con}_T$. Innymi słowy: teoria spełniająca te założenia nie dowodzi własnej niesprzeczności w przyjętej standardowej formalizacji.
\end{twierdzenie}
\begin{proof}
Wybieramy to samo $G$ co w Twierdzeniu~\ref{res:9:2}; z~\eqref{eq:9:1} mamy $T\vdash G\leftrightarrow\neg\Box G$.

\emph{Krok 1: od dowodu $G$ do dowodu sprzeczności.} Z kierunku $G\to\neg\Box G$, przez (D1), a następnie (D2), uzyskujemy
\begin{equation}\label{eq:12:1}
T\vdash\Box G\to\Box\neg\Box G.
\end{equation}
Z (D3) mamy równocześnie
\begin{equation}\label{eq:12:2}
T\vdash\Box G\to\Box\Box G.
\end{equation}
Tautologia $\neg\Box G\to(\Box G\to\bot)$ jest twierdzeniem $T$. Rozpiszmy dokładnie jej użycie, aby nie ukryć zagnieżdżeń $\Box$:
\begin{align*}
T&\vdash\Box(\neg\Box G\to(\Box G\to\bot)) &&\text{(D1)},\\
T&\vdash\Box\neg\Box G\to\Box(\Box G\to\bot) &&\text{(D2)},\\
T&\vdash\Box(\Box G\to\bot)\to(\Box\Box G\to\Box\bot) &&\text{(D2)}.
\end{align*}
Składając dwie ostatnie implikacje, otrzymujemy $T\vdash\Box\neg\Box G\to(\Box\Box G\to\Box\bot)$. Po połączeniu z~\eqref{eq:12:1} i~\eqref{eq:12:2}:
\begin{equation}\label{eq:12:3}
T\vdash\Box G\to\Box\bot.
\end{equation}
Nie użyliśmy zasady $\Box\theta\to\theta$; właśnie jej niedopuszczalność jest istotna dla całego argumentu.

\emph{Krok 2: niesprzeczność implikuje zdanie Gödla wewnątrz teorii.} Kontrapozycja~\eqref{eq:12:3} daje $T\vdash\neg\Box\bot\to\neg\Box G$. Ponieważ $\neg\Box\bot$ to $\operatorname{Con}_T$, a $\neg\Box G\to G$ wynika z punktu stałego, otrzymujemy
\begin{equation}\label{eq:12:4}
T\vdash\operatorname{Con}_T\to G.
\end{equation}
\emph{Krok 3: przejście do metateorii.} Gdyby $T\vdash\operatorname{Con}_T$, przez modus ponens w~\eqref{eq:12:4} otrzymalibyśmy $T\vdash G$. Tymczasem niesprzeczność $T$ daje $T\nvdash G$ na mocy punktu (a) Twierdzenia~\ref{res:9:2}. Sprzeczność kończy dowód. Nie potrzeba tu $\omega$-niesprzeczności.
\end{proof}

\begin{wniosek}[Nieformalny sens]\label{res:12:3}
Jeśli PA jest niesprzeczna, to $\mathrm{PA}\nvdash\operatorname{Con}_{\mathrm{PA}}$. Dowód niesprzeczności w tej formalizacji wymaga więc środków niewyprowadzalnych w samej PA. Klasycznym przykładem jest dowód Gentzena wykorzystujący odpowiednią zasadę indukcji pozaskończonej związaną z $\varepsilon_0$; tego dodatkowego twierdzenia nie dowodzimy tutaj.
\end{wniosek}

